\chapter{Исследовательский раздел}

В данном разделе представляются постановка и проведение исследования разработанного метода.

\section{Постановка исследования}

Задачами проводимого исследования являются:
\begin{itemize}
	\item определение качества разработанного метода на наборе данных FaceForensics;
	\item сравнение разработанного метода с рассмотренными решениями;
	\item классификация разработанного метода.
\end{itemize}

\section{Проведение исследования}

\subsection{Качество разработанного метода на наборе данных FaceForensics}
\label{sec:ff-results}

Качество разработанного метода на наборе данных FaceForensics было определено с использованием множества образцов, содержащего только изображения, созданные методом Face2Face. Данное множество образцов содержит 480 изображений только из набора данных FaceForensics и является подмножеством валидационной части обучающей выборки, описанной в разделе~\ref{sec:training}.

Результаты определения качества разработанного метода представлены в таблицах~\ref{tab:ff-results-1}~и~\ref{tab:ff-results-2}. По таблице~\ref{tab:ff-results-1} была построены матрица ошибок, представленная на рисунке~\ref{img:ff-error-matrix}.

\begin{table}[h]
	\centering
	\caption{Значения количества истинных и ложных срабатываний разработанного метода при определении качества на наборе данных FaceForensics}
	\label{tab:ff-results-1}
	\begin{tabularx}{\textwidth}{|Y|Y|Y|Y|}
		\hline
		\textbf{Количество истинно положительных срабатываний, шт.} & \textbf{Количество истинно отрицательных срабатываний, шт.} & \textbf{Количество ложно положительных срабатываний, шт.} & \textbf{Количество ложно отрицательных срабатываний, шт.} \\
		\hline
		238 & 236 & 2 & 4 \\
		\hline
	\end{tabularx}
\end{table}

\begin{table}[h]
	\centering
	\caption{Результаты определения качества разработанного метода на наборе данных FaceForensics}
	\label{tab:ff-results-2}
	\begin{tabularx}{\textwidth}{|Y|Y|Y|Y|}
		\hline
		\textbf{Точность, \%} & \textbf{Полнота, \%} & \textbf{Сходимость, \%} & \textbf{F-мера, \%} \\
		\hline
		98.75 & 98.35 & 99.17 & 98.76 \\
		\hline
	\end{tabularx}
\end{table}

\includeimage
{ff-error-matrix}
{f} % Обтекание (без обтекания)
{H}
{\textwidth} % Ширина рисунка
{Матрица ошибок при определении качества разработанного метода на наборе данных FaceForensics}

По полученным результатам оценки качества разработанного метода на наборе данных FaceForensics можно сделать вывод, что метод отвечает требованию к точности, предъявленному в разделе~\ref{sec:method-requirements}, имея значение точности 98.75\%, и демонстрирует относительно высокую обобщающую способность на поддельных изображениях, полученных методом Face2Face.

\subsection{Сравнение разработанного метода с рассмотренными решениями}

Сравнение разработанного метода с рассмотренными решениями было проведено по критериям, представленным в таблице~\ref{tab:compare-criteria}. Результаты сравнения  разработанного метода с рассмотренными решениями представлены в таблице~\ref{tab:methods-compare-final}.

\begin{table}[h]
	\centering
	\caption{Сравнение разработанного метода с рассмотренными решениями}
	\label{tab:methods-compare-final}
	\begin{tabularx}{\textwidth}{|>{\centering\arraybackslash}p{3.4cm}|Y|Y|>{\centering\arraybackslash}p{3.5cm}|}
		\hline
		\textbf{Метод} & 
		\textbf{Точность на FaceForensics, \%} &
		\textbf{Количество параметров, шт.} &
		\textbf{Пороговое значение PSNR, дБ.} \\
		\hline
		Meso-4 & 94.6 & 27\,997 & 30 \\
		\hline
		Meso Inception-4 & 96.8 & 28\,615 & 28 \\
		\hline
		Capsule Forensics & 99.37 & 2\,800\,000 & 34 \\
		\hline
		Vision Transformer & 97.0 & 84\,000\,000 & 24 \\
		\hline
		Разработанный метод & 98.75 & 91\,300\,000 & 24 \\
		\hline
	\end{tabularx}
\end{table}

\includeimage
{ff-compare-histogram}
{f}
{H}
{\textwidth}
{Гистограмма значений точности разработанного метода и Capsule Forensics на сжатых изображениях набора данных FaceForensics~\cite{capsforensics}}

\newpage

Из полученных результатов можно сделать вывод, что разработанный метод обнаружения признаков генерации на изображении уступает Capsule Forensics по точности на незашумленных изображениях FaceForensics, однако превосходит остальные рассмотренные свёрточные и трансформерные архитектуры, при этом характеризуется наибольшим количеством параметров среди всех сравниваемых решений.

Сравнение разработанного метода с Capsule Forensics было произведено на множестве данных, описанном в разделе~\ref{sec:ff-results}, предварительно обработанном алгоритмом сжатия JPEG~\cite{jpeg} с параметром качества равным 23. Результаты сравнения точности разработанного метода с Capsule Forensics представлены на гистограмме значений точности, изображенной на рисунке~\ref{img:ff-compare-histogram}.

Из полученных результатов можно сделать вывод, что разработанный метод превосходит Capsule Forensics на низкокачественных изображениях, обладая значением точности равным 97.93\%, то есть имеет более высокую устойчивость к шумам.

\subsection{Классификация разработанного метода}

Классификация разработанного метода была проведена по критериям, представленным в таблице~\ref{tab:classification-criteria}. Результаты классификации разработанного метода представлены в таблице~\ref{tab:methods-classification-final}.

\begin{table}[h]
	\centering
	\caption{Классификация методов обнаружения признаков генерации на изображении}
	\label{tab:methods-classification-final}
	\begin{tabularx}{\textwidth}{|>{\centering\arraybackslash}p{3.4cm}|Y|>{\centering\arraybackslash}p{3.4cm}|Y|}
		\hline
		\textbf{Метод} & \textbf{Уровень декомпозиции изображения} & \textbf{Вычисли-тельная сложность} & \textbf{Устойчивость к шумам} \\
		\hline
		Meso-4 & Мезоскопические & Компактные & Неустойчивые \\
		\hline
		Meso Inception-4 & Мезоскопические & Компактные & Устойчивые \\
		\hline
		Capsule Forensics & Макроскопические & Средние & Неустойчивые \\
		\hline
		Vision Transformer & Холистические & Крупные & Устойчивые \\
		\hline
		Разработанный метод & Холистические & Крупные & Устойчивые \\
		\hline
	\end{tabularx}
\end{table}

\newpage

Из полученных результатов можно сделать вывод, что разработанный метод относится к классу <<устойчивых>> методов, что соответствует предъявленным требованиям.

\section{Выводы из исследовательского раздела}

В результате исследования разработанного метода обнаружения признаков генерации на изображении с использованием нейронной сети было установлено:
\begin{itemize}
	\item разработанный метод превосходит по точности на наборе данных FaceForensics все рассмотренные решения за исключением Capsule Forensics, обладая значением точности равным 98.67\%;
	\item разработанный метод превосходит по точности Capsule Forensics на низкокачественных изображениях набора данных FaceForensics, обладая значением точности равным 97.93\%;
	\item разработанный метод относится к классу <<устойчивых>> методов и содержит не более $10^8$ параметров.
\end{itemize}

Из полученных результатов следует, что разработанный метод отвечает всем предъявленным требованиям, сформированным в разделе~\ref{sec:method-requirements}.